% --- Archon physics formalization source begin ---
% archon:physics
% archon:covers PhyXMiniProblems/problem_phyx_mini_0058.lean
% archon:source-report reports/phyx_mini/problem_phyx_mini_0058.source.json

\chapter{Physics problem phyx\_mini\_0058}
\label{ch:PhyXMiniProblems_problem_phyx_mini_0058}

\paragraph{Problem source.}
\#\# Physical scenario

A stamp collector uses a magnifying lens that sits 2.0 cm above the stamp. The magnification is 4.0.

\#\# Question

What is the focal length of the lens?

\#\# Answer choices

- A: A: 1.4cm

- B: B: 1.6cm

- C: C: 2.7cm

- D: D: 2.4cm

\#\# Figure caption (auxiliary; use the image as primary evidence)

This image illustrates the process of image formation by a lens. Key elements include:

1. **Lens**: Represented by a blue double convex shape in the upper part of the image, designated as the "Lens plane."

2. **Focal Point**: Indicated on the lens axis, labeled as "Focal point."

3. **Ray Diagram**: Purple lines show the path of light rays:
   - Solid lines represent actual light paths.
   - Dashed lines indicate the extension of these rays to form the virtual image.

4. **Object**: A horizontal bar labeled “Stamp,” positioned below the lens plane.

5. **Image Formation**: The virtual image location is marked “Virtual image,” below the object and lens.

6. **Measurements**:
   - Distance \textbackslash{}( s = 2.0 \textbackslash{}, \textbackslash{}text\{cm\} \textbackslash{}) from the lens to the object.
   - Image distance \textbackslash{}( s' = -4.0s \textbackslash{}), indicating the virtual nature of the image (negative value).

This is a typical setup used in optics to illustrate the principles of lens magnification and virtual image creation.

\#\# Dataset metadata

Geometrical Optics; Physical Model Grounding Reasoning, Multi-Formula Reasoning

\paragraph{Recorded answer/context.}
C: C: 2.7cm

\paragraph{Figure/image path.}
/root/proposal\_for\_physic/science-mango/phyx\_mini\_run/phyx\_data/test\_image/58.png

\paragraph{Formalization target.}
create a compiling Lean file with sorry bodies at `PhyXMiniProblems/problem\_phyx\_mini\_0058.lean`. The Lean declarations must preserve the physical quantities, dimensions or dimensional roles, figure labels, governing-law hypotheses, and final relation expressed by this problem.
Use Mathlib/Physlib names found through LeanExplore where available. If a physics API is missing, introduce faithful local abstractions rather than scalar placeholder aliases.

\begin{theorem}[Physics formalization target]
\label{thm:physics:phyx_mini_0058:target}
\lean{PhyXMiniProblems.ProblemPhyXMini0058.problem_phyx_mini_0058}
\uses{def:physics:phyx-mini-0058:phyxminiproblems-problemphyxmini0058-lengthincentimeters, def:physics:phyx-mini-0058:phyxminiproblems-problemphyxmini0058-stampmagnifiersetup, def:physics:phyx-mini-0058:phyxminiproblems-problemphyxmini0058-matchesproblemreadouts, def:physics:phyx-mini-0058:phyxminiproblems-problemphyxmini0058-matchesmagnifierfigure, def:physics:phyx-mini-0058:phyxminiproblems-problemphyxmini0058-hasphysicalopticalbranch, def:physics:phyx-mini-0058:phyxminiproblems-problemphyxmini0058-satisfiesparaxialthinlenslaws, def:physics:phyx-mini-0058:phyxminiproblems-problemphyxmini0058-roundstonearesttenthcentimeter, def:physics:phyx-mini-0058:phyxminiproblems-problemphyxmini0058-isclosestfocallengthanswer}
The assigned autoformalize agent should translate this physics problem into Lean declarations in the covered file, with theorem and lemma proof bodies written as `by sorry`.
\end{theorem}
\begin{proof}
This is an autoformalization task, not a proof task. Produce statements that can later be proved by the physics prover without weakening the physical model.
\end{proof}

% --- Archon declaration topology begin ---
\section{Declaration topology}
\begin{definition}[Length Quantity]
  \label{def:physics:phyx-mini-0058:phyxminiproblems-problemphyxmini0058-lengthquantity}
  \lean{PhyXMiniProblems.ProblemPhyXMini0058.LengthQuantity}
  A signed physical length represented coherently in every choice of units
\end{definition}
\begin{proof}
  This is the declared model definition of Length Quantity.
\end{proof}
\begin{definition}[centimeter Unit Choices]
  \label{def:physics:phyx-mini-0058:phyxminiproblems-problemphyxmini0058-centimeterunitchoices}
  \lean{PhyXMiniProblems.ProblemPhyXMini0058.centimeterUnitChoices}
  SI base units with centimeters selected as the length unit
\end{definition}
\begin{proof}
  This is the declared model definition of centimeter Unit Choices.
\end{proof}
\begin{definition}[length In Centimeters]
  \label{def:physics:phyx-mini-0058:phyxminiproblems-problemphyxmini0058-lengthincentimeters}
  \lean{PhyXMiniProblems.ProblemPhyXMini0058.lengthInCentimeters}
  \uses{def:physics:phyx-mini-0058:phyxminiproblems-problemphyxmini0058-lengthquantity, def:physics:phyx-mini-0058:phyxminiproblems-problemphyxmini0058-centimeterunitchoices}
  The signed scalar centimeter readout of a physical length
\end{definition}
\begin{proof}
  This is the declared model definition of length In Centimeters.
\end{proof}
\begin{definition}[Figure Location]
  \label{def:physics:phyx-mini-0058:phyxminiproblems-problemphyxmini0058-figurelocation}
  \lean{PhyXMiniProblems.ProblemPhyXMini0058.FigureLocation}
  Labels attached to locations in the magnifying-lens figure
\end{definition}
\begin{proof}
  This is the declared model definition of Figure Location.
\end{proof}
\begin{definition}[Ray Style]
  \label{def:physics:phyx-mini-0058:phyxminiproblems-problemphyxmini0058-raystyle}
  \lean{PhyXMiniProblems.ProblemPhyXMini0058.RayStyle}
  Whether a pictured ray is an actual solid ray or a dashed backward extension
\end{definition}
\begin{proof}
  This is the declared model definition of Ray Style.
\end{proof}
\begin{definition}[Thin Lens Kind]
  \label{def:physics:phyx-mini-0058:phyxminiproblems-problemphyxmini0058-thinlenskind}
  \lean{PhyXMiniProblems.ProblemPhyXMini0058.ThinLensKind}
  Optical action of the thin lens
\end{definition}
\begin{proof}
  This is the declared model definition of Thin Lens Kind.
\end{proof}
\begin{definition}[Lens Profile]
  \label{def:physics:phyx-mini-0058:phyxminiproblems-problemphyxmini0058-lensprofile}
  \lean{PhyXMiniProblems.ProblemPhyXMini0058.LensProfile}
  Geometric profile drawn for the lens
\end{definition}
\begin{proof}
  This is the declared model definition of Lens Profile.
\end{proof}
\begin{definition}[Stamp Magnifier Setup]
  \label{def:physics:phyx-mini-0058:phyxminiproblems-problemphyxmini0058-stampmagnifiersetup}
  \lean{PhyXMiniProblems.ProblemPhyXMini0058.StampMagnifierSetup}
  \uses{def:physics:phyx-mini-0058:phyxminiproblems-problemphyxmini0058-lengthquantity, def:physics:phyx-mini-0058:phyxminiproblems-problemphyxmini0058-figurelocation, def:physics:phyx-mini-0058:phyxminiproblems-problemphyxmini0058-raystyle, def:physics:phyx-mini-0058:phyxminiproblems-problemphyxmini0058-thinlenskind, def:physics:phyx-mini-0058:phyxminiproblems-problemphyxmini0058-lensprofile}
  The dimensionful quantities and labeled geometry of the stamp magnifier. objectDistance is the positive figure label s; signedImageDistance is the signed label s'; and focalLength is the positive distance f from the lens plane to the displayed focal point. transverseMagnification is signed and dimensionless, so an upright virtual image has positive magnification
\end{definition}
\begin{proof}
  This is the declared model definition of Stamp Magnifier Setup.
\end{proof}
\begin{definition}[Matches Problem Readouts]
  \label{def:physics:phyx-mini-0058:phyxminiproblems-problemphyxmini0058-matchesproblemreadouts}
  \lean{PhyXMiniProblems.ProblemPhyXMini0058.MatchesProblemReadouts}
  \uses{def:physics:phyx-mini-0058:phyxminiproblems-problemphyxmini0058-lengthincentimeters, def:physics:phyx-mini-0058:phyxminiproblems-problemphyxmini0058-stampmagnifiersetup}
  Numerical information in the problem statement: the lens is 2.0 cm above the stamp and its signed transverse magnification is +4.0. No focal-length value occurs in these data
\end{definition}
\begin{proof}
  This is the declared model definition of Matches Problem Readouts.
\end{proof}
\begin{definition}[Matches Magnifier Figure]
  \label{def:physics:phyx-mini-0058:phyxminiproblems-problemphyxmini0058-matchesmagnifierfigure}
  \lean{PhyXMiniProblems.ProblemPhyXMini0058.MatchesMagnifierFigure}
  \uses{def:physics:phyx-mini-0058:phyxminiproblems-problemphyxmini0058-lengthincentimeters, def:physics:phyx-mini-0058:phyxminiproblems-problemphyxmini0058-stampmagnifiersetup}
  Information read directly from the primary figure. Coordinates increase from the stamp side through the lens toward the outgoing side. Thus the virtual image lies farther on the incoming side than the stamp, while the displayed focal point lies beyond the lens plane. The annotation s' = -4.0 s is kept as a figure readout rather than treated as the requested focal-length answer
\end{definition}
\begin{proof}
  This is the declared model definition of Matches Magnifier Figure.
\end{proof}
\begin{definition}[Has Physical Optical Branch]
  \label{def:physics:phyx-mini-0058:phyxminiproblems-problemphyxmini0058-hasphysicalopticalbranch}
  \lean{PhyXMiniProblems.ProblemPhyXMini0058.HasPhysicalOpticalBranch}
  \uses{def:physics:phyx-mini-0058:phyxminiproblems-problemphyxmini0058-lengthincentimeters, def:physics:phyx-mini-0058:phyxminiproblems-problemphyxmini0058-stampmagnifiersetup}
  The real-object, virtual-image, magnifying branch depicted in the figure
\end{definition}
\begin{proof}
  This is the declared model definition of Has Physical Optical Branch.
\end{proof}
\begin{definition}[Satisfies Paraxial Thin Lens Laws]
  \label{def:physics:phyx-mini-0058:phyxminiproblems-problemphyxmini0058-satisfiesparaxialthinlenslaws}
  \lean{PhyXMiniProblems.ProblemPhyXMini0058.SatisfiesParaxialThinLensLaws}
  \uses{def:physics:phyx-mini-0058:phyxminiproblems-problemphyxmini0058-stampmagnifiersetup}
  The paraxial governing laws for the setup, stated in every choice of units. The first field is the denominator-free form of 1 / f = 1 / s + 1 / s'. The second is the denominator-free signed magnification law m = -s' / s. These fields are general physical laws and contain no numerical focal-length answer
\end{definition}
\begin{proof}
  This is the declared model definition of Satisfies Paraxial Thin Lens Laws.
\end{proof}
\begin{definition}[Answer Choice]
  \label{def:physics:phyx-mini-0058:phyxminiproblems-problemphyxmini0058-answerchoice}
  \lean{PhyXMiniProblems.ProblemPhyXMini0058.AnswerChoice}
  Labels of the four displayed multiple-choice answers
\end{definition}
\begin{proof}
  This is the declared model definition of Answer Choice.
\end{proof}
\begin{definition}[answer Focal Length In Centimeters]
  \label{def:physics:phyx-mini-0058:phyxminiproblems-problemphyxmini0058-answerfocallengthincentimeters}
  \lean{PhyXMiniProblems.ProblemPhyXMini0058.answerFocalLengthInCentimeters}
  \uses{def:physics:phyx-mini-0058:phyxminiproblems-problemphyxmini0058-answerchoice}
  The focal-length readout printed beside each answer choice, in centimeters
\end{definition}
\begin{proof}
  This is the declared model definition of answer Focal Length In Centimeters.
\end{proof}
\begin{definition}[Rounds To Nearest Tenth Centimeter]
  \label{def:physics:phyx-mini-0058:phyxminiproblems-problemphyxmini0058-roundstonearesttenthcentimeter}
  \lean{PhyXMiniProblems.ProblemPhyXMini0058.RoundsToNearestTenthCentimeter}
  \uses{def:physics:phyx-mini-0058:phyxminiproblems-problemphyxmini0058-lengthquantity, def:physics:phyx-mini-0058:phyxminiproblems-problemphyxmini0058-lengthincentimeters, def:physics:phyx-mini-0058:phyxminiproblems-problemphyxmini0058-answerchoice, def:physics:phyx-mini-0058:phyxminiproblems-problemphyxmini0058-answerfocallengthincentimeters}
  Agreement with a displayed focal length rounded to the nearest tenth centimeter
\end{definition}
\begin{proof}
  This is the declared model definition of Rounds To Nearest Tenth Centimeter.
\end{proof}
\begin{definition}[Is Closest Focal Length Answer]
  \label{def:physics:phyx-mini-0058:phyxminiproblems-problemphyxmini0058-isclosestfocallengthanswer}
  \lean{PhyXMiniProblems.ProblemPhyXMini0058.IsClosestFocalLengthAnswer}
  \uses{def:physics:phyx-mini-0058:phyxminiproblems-problemphyxmini0058-lengthquantity, def:physics:phyx-mini-0058:phyxminiproblems-problemphyxmini0058-lengthincentimeters, def:physics:phyx-mini-0058:phyxminiproblems-problemphyxmini0058-answerchoice, def:physics:phyx-mini-0058:phyxminiproblems-problemphyxmini0058-answerfocallengthincentimeters}
  A selected option is strictly closer to the exact focal length than every rival
\end{definition}
\begin{proof}
  This is the declared model definition of Is Closest Focal Length Answer.
\end{proof}
\begin{lemma}[signed Image Distance eq neg eight]
  \label{lem:physics:phyx-mini-0058:phyxminiproblems-problemphyxmini0058-signedimagedistance-eq-neg-eight}
  \lean{PhyXMiniProblems.ProblemPhyXMini0058.signedImageDistance_eq_neg_eight}
  \uses{def:physics:phyx-mini-0058:phyxminiproblems-problemphyxmini0058-lengthincentimeters, def:physics:phyx-mini-0058:phyxminiproblems-problemphyxmini0058-stampmagnifiersetup, def:physics:phyx-mini-0058:phyxminiproblems-problemphyxmini0058-matchesproblemreadouts, def:physics:phyx-mini-0058:phyxminiproblems-problemphyxmini0058-satisfiesparaxialthinlenslaws}
  The signed magnification law and the two problem readouts determine the virtual-image distance as s' = -8 cm
\end{lemma}
\begin{proof}
  The conclusion follows from the governing relations and figure readouts named in the statement of signed Image Distance eq neg eight.
\end{proof}
\begin{lemma}[focal Length eq eight thirds]
  \label{lem:physics:phyx-mini-0058:phyxminiproblems-problemphyxmini0058-focallength-eq-eight-thirds}
  \lean{PhyXMiniProblems.ProblemPhyXMini0058.focalLength_eq_eight_thirds}
  \uses{def:physics:phyx-mini-0058:phyxminiproblems-problemphyxmini0058-lengthincentimeters, def:physics:phyx-mini-0058:phyxminiproblems-problemphyxmini0058-stampmagnifiersetup, def:physics:phyx-mini-0058:phyxminiproblems-problemphyxmini0058-matchesproblemreadouts, def:physics:phyx-mini-0058:phyxminiproblems-problemphyxmini0058-satisfiesparaxialthinlenslaws}
  The Gaussian thin-lens equation then determines the exact focal length f = 8/3 cm
\end{lemma}
\begin{proof}
  The conclusion follows from the governing relations and figure readouts named in the statement of focal Length eq eight thirds.
\end{proof}
% --- Archon declaration topology end ---
% --- Archon physics formalization source end ---
